\chapterhead{1}{Energy You Can Lift}{A drop and a bounce, written as heights}\label{ch1}

Lift a ball and hold it up. It is not moving, yet it can do something the moment you let go: it falls, it speeds up and it hits the floor. We say a lifted ball has \textbf{stored energy}, and the higher it is lifted the more it has.

Energy is hard to measure at home. So we use a stand-in. For one ball, the height of the lift tells us how much stored energy it has, and \textbf{1~unit is 1~cm of lift}. A ball held at 40~cm has 40 units. This works only for the \emph{same ball}.

\begin{center}
\sffamily\small\renewcommand{\arraystretch}{1.5}
\begin{tabular}{@{}p{3.4cm}p{3.7cm}p{3.7cm}@{}}
\toprule
\textbf{Model} & \textbf{What it keeps} & \textbf{What it drops}\\
\midrule
Height as energy & how much the lift stored & the ball's weight and shape\\
A bounce height & how much came back up & where the rest went\\
An energy unit & a number we can add & what energy really is\\
\bottomrule
\end{tabular}
\end{center}

\begin{tryit}{Drop and bounce}
\textbf{You need:} a bouncy ball, a ruler or tape measure stuck upright to a wall, a hard floor.
\begin{enumerate}[leftmargin=1.6em]
\item Hold the bottom of the ball at 40~cm. Drop it. Do not throw it.
\item Watch how high the bottom of the ball gets on its \emph{first} bounce. Ask a helper to watch too.
\item Repeat five times and write each bounce height. Then work out the average.
\end{enumerate}
\par\medskip
\renewcommand{\arraystretch}{1.6}
\begin{tabular}{|l|c|c|c|c|c|c|}
\hline
 & \textbf{1} & \textbf{2} & \textbf{3} & \textbf{4} & \textbf{5} & \textbf{Average}\\ \hline
\textbf{Bounce (cm)} & & & & & & \\ \hline
\end{tabular}
\end{tryit}
\adult{Drop the ball from low heights only and keep clear of windows, stairs and lamps. Small balls can be swallowed by little children and pets.}

\begin{worked}{Practice numbers, not from a real test}
Bounces: 29, 27, 28, 28, 28. Total 140, so the average is $140\div5=28$~cm.\\
Start: 40 units. Back up: 28 units. The gap is $40-28=12$ units.\\
Those 12 units did not vanish. We cannot measure them, but warmth in the ball and floor and the sound of the thud are the likely places.
\end{worked}

\begin{center}
\begin{tikzpicture}
  \node[bigcir,minimum size=2.4cm] (a) at (0,0) {}; \node[font=\sffamily\scriptsize,align=center] at (0,0) {lift\\40};
  \node[cir] (b) at (4.6,0.9) {28}; \node[cir] (c) at (4.6,-0.9) {12};
  \draw[popA] (a) -- node[above,note] {bounce} (b);
  \draw[popA] (a) -- node[below,note] {warmth, sound} (c);
  \node[note,anchor=west,text width=2.6cm,align=left] at (5.6,0) {$28+12=40$};
\end{tikzpicture}
\end{center}

\begin{multix}[hmm]\multixbot[hmm]\tcblower
The ball came back up to 28, not 40, so 12 units were destroyed. I have a very tidy way of saying it. Is it right?
\end{multix}

\trust{We dropped from the same height five times, with the same ball on the same floor, and averaged the bounces.}
{Every bounce was within 1 or 2~cm of the others, so the average is steady.}
{Where the other 12 units went. Warmth and sound are the likeliest places, but our kit cannot measure them.}

\begin{keepbox}
We turned a bounce into two heights. \textbf{We kept} \wline[3.4cm] \textbf{and we set aside} \wline[3.4cm].
\end{keepbox}

\begin{tryit}{Three heights}
Drop the same ball from 20~cm, 40~cm and 60~cm. Take three bounces at each height and write the average.\par\medskip
\renewcommand{\arraystretch}{1.6}
\begin{tabular}{|l|c|c|c|}
\hline
\textbf{Drop (units)} & \textbf{Average bounce (cm)} & \textbf{Lost (units)} & \textbf{Bounce $\div$ drop}\\ \hline
20 & & & \\ \hline
40 & & & \\ \hline
60 & & & \\ \hline
\end{tabular}
\end{tryit}

\begin{worked}{Practice numbers}
Drop 20: bounce 14, lost 6. Drop 40: bounce 28, lost 12. Drop 60: bounce 42, lost 18.\\
$14\div20=0.7$, $28\div40=0.7$ and $42\div60=0.7$. This ball gives back about 7 tenths of its lift. Three tenths go elsewhere.
\end{worked}

\clearpage

\begin{puzzle}{Lost}
A ball is dropped from 50 and bounces to 35. How many units are lost? \wline[1.6cm]
\end{puzzle}

\begin{puzzle}{Average}
Five bounces were 31, 29, 30, 30 and 30. Average \wline[1.6cm]~cm. Lost from a 40 drop \wline[1.6cm]
\end{puzzle}

\begin{puzzle}{Which form?}
Write S for stored, M for motion or W for warmth.\\
(a) ball held up high \wline[1cm] \quad (b) ball falling \wline[1cm] \quad (c) floor slightly warm after a hit \wline[1cm]\\
(d) book on a top shelf \wline[1cm] \quad (e) rolling skateboard \wline[1cm]
\end{puzzle}

\begin{puzzle}{Unfair test}
Ana dropped a tennis ball from 40 and Ben dropped a foam ball from 60. They say foam bounces worse. What is wrong with their test? \wline[6cm]
\end{puzzle}

\begin{puzzle}{Next bounces}
A ball bounces to 28, then 20, then 14. How much was lost between the 1st and 2nd bounce? \wline[1.6cm] Total lost from a 40 drop to the 3rd bounce? \wline[1.6cm]
\end{puzzle}

\clearpage

\begin{puzzle}{Fill the row}
Drop 80, bounce 56. Lost \wline[1.6cm] units.
\end{puzzle}

\begin{puzzle}{Which sentence?}
(a) ``12 units were destroyed.'' (b) ``12 units went somewhere we did not measure.'' Which is better supported? \wline[1.6cm]
\end{puzzle}

\begin{puzzle}{Draw the circles}
A ball starts with 40 and bounces to 28. Fill in the empty circle.
\begin{center}
\begin{tikzpicture}
  \node[bigcir,minimum size=2.4cm] (a) at (0,0) {}; \node[font=\sffamily\scriptsize] at (0,0) {lift 40};
  \node[cir] (b) at (4,0.8) {28}; \node[cir,dashed] (c) at (4,-0.8) {};
  \draw[popA] (a) -- (b); \draw[popA] (a) -- (c);
\end{tikzpicture}
\end{center}
\end{puzzle}

\begin{puzzle}{Wrap}
A student wrote ``The ball lost energy.'' Wrap it with a clearer claim. \wline[6cm]
\end{puzzle}

\begin{puzzle}{Design a drop}
Choose a different ball or floor and say what you would keep the same. \wline[6cm]
\end{puzzle}
\clearpage

\thispagestyle{fancy}
\subhead{Bounce record}
Try three balls or three floors. Keep the drop height the same each time.\par\medskip
\renewcommand{\arraystretch}{3.0}
\begin{tabularx}{\linewidth}{|>{\sffamily\small\bfseries}p{2.6cm}|X|X|X|X|}
\hline
Ball and floor & Drop (units) & Average bounce & Lost & Bounce $\div$ drop \\ \hline
 & & & & \\ \hline
 & & & & \\ \hline
 & & & & \\ \hline
\end{tabularx}

\vspace{8pt}
Which combination gave back the most? \wline[7cm]\\[10pt]
Which lost the most, and what might it have lost it to? \wline[5cm]\\[10pt]
Sketch your best bounce and mark the heights.\par\medskip
\begin{tikzpicture}\draw[ink!70,line width=1pt,rounded corners=3pt] (0,0) rectangle (11.4,5.2);\end{tikzpicture}
\clearpage
